% =====================================================================
%  Bruno Le Floch, "A slow review of the AGT correspondence"
%  arXiv:2006.14025v3 (JPhysA review, Jan 2022) —— 中文译注本
%
%  编译：LuaLaTeX（本机 xdvipdfmx 不稳定，统一用 lualatex，编译 3 遍）
%  用法：lualatex -interaction=nonstopmode main.tex  （×3）
%
%  约定：
%   * 公式编号用 \tag{} 复刻原文编号（如 (3.13)），与原文一一对应；
%   * 图 1--5 为原文矢量图按 400 dpi 截取，图注为译者所译；
%   * % [Tutor Note: ...] 为译者批注，不进入排版，汇总见 译注汇总.md；
%   * 原文每节开头的灰色「要点框」用 keybox 环境排版（标题为原文粗体标题）。
% =====================================================================
\documentclass[11pt,a4paper]{article}

\usepackage[UTF8]{ctex}
\usepackage{amsmath,amssymb,amsthm,mathtools}
\usepackage{bm}
\usepackage{graphicx}
\usepackage[a4paper,left=2.5cm,right=2.5cm,top=2.6cm,bottom=2.7cm]{geometry}
\usepackage[dvipsnames]{xcolor}
\usepackage{enumitem}
\usepackage{booktabs}
\usepackage{array}
\usepackage{multirow}
\usepackage{longtable}
\usepackage{multicol}
\usepackage{fancyhdr}
\usepackage{caption}
\usepackage{tcolorbox}
\usepackage{tikz}
\usepackage[hidelinks,unicode]{hyperref}
\hypersetup{pdftitle={AGT 对应关系慢讲（中译注本）},
  pdfauthor={Bruno Le Floch; 中译/批注: 本译本}}

\captionsetup{font=small,labelsep=quad}
\setlength{\emergencystretch}{3em}
\allowdisplaybreaks[1]
\graphicspath{{figures/}{../figures/}}

% ---------- 篇标题：原文用 “Part I/II/III”，中文版由译者直接写“第一篇/第二篇/第三篇”，
%            故抑制 LaTeX 自动的 “Part I” 前缀（否则会重复出现）----------
\renewcommand{\partname}{}
\renewcommand{\thepart}{}

% ---------- 页眉 ----------
\pagestyle{fancy}
\fancyhf{}
\fancyhead[L]{\small\itshape A slow review of the AGT correspondence}
\fancyhead[R]{\small\thepage}
\renewcommand{\headrulewidth}{0.4pt}

% ---------- 定理类环境（编号一律由原文手工给出） ----------
\theoremstyle{plain}
\newtheorem*{theorem*}{定理}
\newtheorem*{proposition*}{命题}
\newtheorem*{lemma*}{引理}
\theoremstyle{definition}
\newtheorem*{definition*}{定义}
\newtheorem*{remark*}{注}

% 原文的「Exercise x.y.」：正文中为粗体标签 + 斜体正文
\newenvironment{exercise}[1]%
  {\par\medskip\noindent\textbf{习题 #1.}\itshape\ }{\par\medskip}

% 原文每节开头的灰色要点框（深色标题条 + 浅灰底）
% 用法：\begin{keybox}{标题文字} \begin{itemize}...\end{itemize} \end{keybox}
\newtcolorbox{keybox}[1]{%
  colback=black!5, colframe=black!55, colbacktitle=black!75, coltitle=white,
  fonttitle=\bfseries\small, title={#1}, boxrule=0.6pt, arc=2pt,
  left=6pt, right=6pt, top=4pt, bottom=4pt, toptitle=2pt, bottomtitle=2pt,
  fontupper=\small}

% 无标题版要点框（原文个别要点框没有深色标题条，如 p.70 与 §10）
\newtcolorbox{keyboxplain}{%
  colback=black!5, colframe=black!55, boxrule=0.6pt, arc=2pt,
  left=6pt, right=6pt, top=4pt, bottom=4pt, fontupper=\small}

% 原文作者用粗体行内小标题（如 "**Tame punctures.**"）——
% 中文译为 \para{...}，独立成段、粗体、不缩进。
\newcommand{\para}[1]{\par\medskip\noindent\textbf{#1}\quad\ignorespaces}

% 原文表格标题（手工编号）
\newcommand{\tabtitle}[2]{\caption*{\small\textbf{表 #1：}#2}}
% 原文图标题（手工编号）
\newcommand{\figtitle}[2]{\caption*{\small\textbf{图 #1：}#2}}

% 参考文献条目（生成脚本按 \refitem{编号}{内容} 输出）
\newcommand{\refitem}[2]{\par\hangindent=1.5em\hangafter=1\noindent [#1]~#2}

% ---------- 常用记号（只能用这里定义过的宏，不要自造 \newcommand） ----------
\newcommand{\R}{\mathbb{R}}
\newcommand{\C}{\mathbb{C}}
\newcommand{\Z}{\mathbb{Z}}
\newcommand{\N}{\mathbb{N}}
\newcommand{\Q}{\mathbb{Q}}
\newcommand{\cN}{\mathcal{N}}
\newcommand{\cO}{\mathcal{O}}
\newcommand{\cM}{\mathcal{M}}
\newcommand{\cB}{\mathcal{B}}
\newcommand{\cC}{\mathcal{C}}
\newcommand{\cD}{\mathcal{D}}
\newcommand{\cE}{\mathcal{E}}
\newcommand{\cF}{\mathcal{F}}
\newcommand{\cG}{\mathcal{G}}
\newcommand{\cH}{\mathcal{H}}
\newcommand{\cI}{\mathcal{I}}
\newcommand{\cL}{\mathcal{L}}
\newcommand{\cP}{\mathcal{P}}
\newcommand{\cQ}{\mathcal{Q}}
\newcommand{\cS}{\mathcal{S}}
\newcommand{\cT}{\mathcal{T}}
\newcommand{\cU}{\mathcal{U}}
\newcommand{\cV}{\mathcal{V}}
\newcommand{\cW}{\mathcal{W}}
\newcommand{\cX}{\mathcal{X}}
\newcommand{\cY}{\mathcal{Y}}
\newcommand{\cZ}{\mathcal{Z}}
\newcommand{\cA}{\mathcal{A}}
\newcommand{\cJ}{\mathcal{J}}
\newcommand{\cK}{\mathcal{K}}
\newcommand{\cR}{\mathcal{R}}
\newcommand{\fg}{\mathfrak{g}}
\newcommand{\fh}{\mathfrak{h}}
\newcommand{\ft}{\mathfrak{t}}
\newcommand{\fu}{\mathfrak{u}}
\newcommand{\fsu}{\mathfrak{su}}
\newcommand{\fso}{\mathfrak{so}}
\newcommand{\fsp}{\mathfrak{sp}}
\newcommand{\dd}{\mathrm{d}}
\newcommand{\ee}{\mathrm{e}}
\newcommand{\ii}{\mathrm{i}}
\newcommand{\Tr}{\operatorname{Tr}}
\newcommand{\tr}{\operatorname{tr}}
\newcommand{\diag}{\operatorname{diag}}
\newcommand{\sgn}{\operatorname{sgn}}
\newcommand{\id}{\mathrm{id}}
\renewcommand{\Re}{\operatorname{Re}}
\renewcommand{\Im}{\operatorname{Im}}
\newcommand{\rk}{\operatorname{rk}}
\newcommand{\rank}{\operatorname{rank}}
\newcommand{\Pf}{\operatorname{Pf}}
\newcommand{\Res}{\operatorname{Res}}
\newcommand{\const}{\mathrm{const}}
\newcommand{\Toda}{\mathrm{Toda}}
\newcommand{\Liouville}{\mathrm{Liouville}}
\newcommand{\bV}{\widehat{V}}
\newcommand{\vol}{\operatorname{vol}}
\newcommand{\Zbb}{\mathbb{Z}}
\newcommand{\half}{\tfrac{1}{2}}
\newcommand{\ie}{\textit{i.e.}}
\newcommand{\eg}{\textit{e.g.}}
\newcommand{\cf}{\textit{cf.}}
% 原文中的 \bar 简写（避免与 \bar 冲突，统一用标准 \bar）
\newcommand{\bQ}{\bar{Q}}
\newcommand{\bq}{\bar{q}}
\newcommand{\bepsilon}{\bar{\epsilon}}
\newcommand{\btheta}{\bar{\theta}}
